% Source: https://artofproblemsolving.com/community/c3h2173064_tricky_maths
% Stable identity: aops-topic-2173064-q02
Use coordinates $A=(0,0)$, $B=(w,0)$, $C=(w,h)$, $D=(0,h)$. Then $E=\left(0,\dfrac {h}{2}\right)$ and $F=\left(0,\dfrac{3h}{4}\right)$.

\begin{asy}
usepackage("amsmath");
size(6.5cm);
defaultpen(fontsize(10pt));
pair A=(0,0),B=(4,0),C=(4,3),D=(0,3),Ep=(0,1.5),F=(0,2.25),G=(-4,0),H=(-4/3,3);
fill(B--C--G--cycle,rgb(0.85,0.9,1)); fill(B--C--H--cycle,rgb(1,0.85,0.85)+opacity(0.6));
draw(A--B--C--D--cycle,linewidth(0.9));
draw(G--A,dashed); draw(H--D,dashed);
draw(C--G,blue); draw(B--H,red);
dot(A);dot(B);dot(C);dot(D);dot(Ep);dot(F);dot(G);dot(H);
label("$A$",A,S);label("$B$",B,SE);label("$C$",C,NE);label("$D$",D,N);label("$E$",Ep,dir(150));label("$F$",F,NE);
label("$G$",G,SW);label("$H$",H,N);
label("$w$",(A+B)/2,S,fontsize(8pt)); label("$h$",(B+C)/2,E,fontsize(8pt));
label("\mbox{}",(min(currentpicture,user=true)+max(currentpicture,user=true))/2,black); // keeps the text after the figure black when LaTeX typesets the labels inline
\end{asy}

\textbf{Point $G$.} Line $CE$ rises $\dfrac {h}{2}$ over a run of $w$; going from $E$ down to height $0$ (another drop of $\dfrac {h}{2}$) takes a further run of $w$ to the left, so $G=(-w,0)$. (Equivalently, triangles $EAG$ and $EDC$ are congruent, so $AG=DC=w$.)

\textbf{Point $H$.} Line $BF$ goes from $B=(w,0)$ to $F=\left(0,\dfrac{3h}{4}\right)$; it reaches height $h$ after a horizontal run of $\dfrac{4}{3}w$ from $B$, so $H=\left(-\dfrac {w}{3},h\right)$. (Equivalently, triangles $FDH$ and $FAB$ are similar with ratio $\dfrac{FD}{FA}=\dfrac{1}{3}$, so $DH=\dfrac {w}{3}$.)

\textbf{Areas.} Triangle $BCG$ has base $BG=2w$ on line $AB$ and height $h$, so $[BCG]=wh$. Triangle $BCH$ has base $CH=w+\dfrac {w}{3}=\dfrac{4w}{3}$ on line $CD$ and height $h$, so $[BCH]=\dfrac{2wh}{3}$. Hence
\[\dfrac{[BCG]}{[BCH]}=\dfrac{wh}{\dfrac{2wh}{3}}=\dfrac{3}{2},\]
so $m=3$, $n=2$ and $10m+10n+mn=30+20+6=\boxed{56}$.
